\paragraph{A head on the tokens reads the closing speed, and by the frozen rule no better than the reader's outputs (row 143, H-V8F-PROOF-143).} Card ADD81 asks whether the closing speed can be read from more than the reader's three outputs: a small recurrent head over the last 10 or 25 frames, either on the reader's frozen embedding or on the pooled token features of the model (\tier{T0}, reads on logged frames, two seeds, all five control lines passed at the first read; the scorer reproduces the cells of row 139 exactly, and a head fitted on shuffled labels reads $+0.0667$). All four heads read the closing speed, and none is better than the estimate from the outputs by the rule of the card, so the name is that the tokens give no gain. The head on the tokens over 25 frames is nonetheless the best read of the closing speed on record here: r $+0.468$ [$+0.331$, $+0.611$] against $+0.384$ for the ridge estimate, a mean absolute error of 1.785 against 1.876\,m/s; the paired difference of the error is $-0.091$ [$-0.225$, $+0.029$], and an interval that reaches $+0.029$ is not a gain. The heads on the frozen embedding and the token head over 10 frames are worse than their cell, two of them at or above the error of a constant read. One of the two fits of the best head ended at its cap while its loss still fell. The card commits a cost card to the PI as the next item, since a crop of the lead region or a finer encode needs camera frames back on the pod. Before that card is written, two things are measured without a new fit (card ADD84): what the saved table of the best head does as the closing input of the rule, with the estimate of another clip as its control, and what a closing read of a planted quality would buy, the label closing speed with calibrated noise at three levels of r, so that the cost card can name the quality a better read has to reach.
